\documentclass{beamer}
\usepackage[utf8]{inputenc}
\usepackage[T1]{fontenc}
\usepackage[french]{babel}
\usepackage{booktabs}
\usepackage{graphicx}

\usetheme{Madrid}

\title[Évaluation EMOTYC]{Évaluation des performances d'EMOTYC sur CyberAggAdo}
\subtitle{Synthèse du rapport de performances}
\author{Équipe d'évaluation}
\date{\today}

\begin{document}

\begin{frame}
    \titlepage
\end{frame}

\begin{frame}{Ce que le modèle a appris}
    \begin{itemize}
        \item \textbf{EMOTYC} : CamemBERT-base fine-tuné (19 sorties binaires : Emo, modes, base/complexe, 12 catégories).
        \item \textbf{Corpus source (TextToKids)} : Textes journalistiques jeunesse (91\%), encyclopédiques, peu de romans.
        \item Entraînement sur des textes édités, non conversationnels.
        \item \textbf{Conséquence directe} : EMOTYC connaît les marqueurs émotionnels de textes propres, mais pas la pragmatique des interactions courtes et bruitées de CyberAggAdo.
    \end{itemize}
\end{frame}

\begin{frame}{Résultats globaux}
    \begin{table}[]
        \resizebox{\textwidth}{!}{%
        \begin{tabular}{llrrr}
        \toprule
        \textbf{Corpus / Configuration} & \textbf{Template} & \textbf{Macro-F1} & \textbf{Micro-F1} & \textbf{Exact match} \\ \midrule
        TextToKids (avec contexte) & \texttt{bca\_spaced\_context} & 0.739 & 0.897 & n.d. \\
        CyberAggAdo (complet) & \texttt{bca\_context} & 0.291 & 0.472 & 0.346 \\
        CyberAggAdo (orchestré) & \texttt{bca\_spaced\_context} & 0.282 & 0.472 & n.d. \\
        Random sample 120 & \texttt{bca\_spaced\_context} & 0.247 & 0.492 & 0.283 \\ \bottomrule
        \end{tabular}%
        }
    \end{table}
    \vspace{0.5cm}
    \begin{itemize}
        \item \textbf{Effondrement robuste des performances} : Le Micro-F1 chute à environ 0.47 en OOD (Out-Of-Domain).
        \item Le problème persiste sur des sous-échantillons contigus et aléatoires.
    \end{itemize}
\end{frame}

\begin{frame}{Chute par famille de labels}
    \begin{table}[]
        \begin{tabular}{lrr}
        \toprule
        \textbf{Groupe} & \textbf{F1 TextToKids} & \textbf{F1 CyberAggAdo} \\ \midrule
        Emo & 0.930 & 0.657 \\
        Types (Base/Complexe) & 0.882 & 0.418 \\
        Modes & 0.881 & 0.266 \\
        Catégories émotionnelles & 0.652 & 0.248 \\ \midrule
        \textbf{Tous les 19 labels} & \textbf{0.739} & \textbf{0.291} \\ \bottomrule
        \end{tabular}
    \end{table}
    \vspace{0.5cm}
    \begin{itemize}
        \item Le modèle repère la charge émotionnelle globale (\texttt{Emo} F1 = 0.657).
        \item \textbf{Problème de granularité :} Incapacité à choisir la bonne catégorie ou le bon mode dans un registre de cyberharcèlement.
    \end{itemize}
\end{frame}

\begin{frame}{Résultats par label (Sélection)}
    \begin{itemize}
        \item \textbf{Colère (F1 = 0.477) :} Label central, mais confondu avec agressivité et menace de façon indistincte.
        \item \textbf{Dégoût (F1 = 0.000) :} 80 occurrences gold, 0 prédiction. Le modèle ne le détecte jamais.
        \item \textbf{Autre (F1 = 0.130) :} Forte sur-prédiction (125 FP) et précision très faible.
        \item \textbf{Modes (Suggérée / Comportementale) :} Incapacité quasi totale à inférer l'émotion situationnelle (F1 = 0.085) et à lier menaces/actions aux comportements (F1 = 0.148).
    \end{itemize}
\end{frame}

\begin{frame}{Distribution des erreurs et Cohérence}
    \textbf{Distribution des erreurs :}
    \begin{itemize}
        \item Seulement \textbf{34.6\%} des phrases sont exactes sur les 19 labels.
        \item \textbf{45.6\%} des lignes accumulent au moins 3 erreurs (moyenne de 2.23 erreurs/ligne).
    \end{itemize}

    \vspace{0.5cm}
    \textbf{Cohérence structurelle (sorties indépendantes) :}
    \begin{itemize}
        \item 24.3\% des lignes présentent au moins une violation logique.
        \item \textbf{Violation principale :} Émotion prédite sans mode associé (15.9\%).
        \item \textit{Diagnostic :} EMOTYC détecte qu'il se passe quelque chose d'émotionnel, mais ne sait pas comment cela s'exprime dans des messages sarcastiques ou insultants.
    \end{itemize}
\end{frame}

\begin{frame}{Pourquoi le modèle se trompe ? (1/2)}
    \begin{itemize}
        \item \textbf{Changement de domaine (Domain Shift) :} Textes édités (source) vs. conversations adolescentes, insultes, abréviations, sarcasme (cible).
        \item \textbf{Agressivité $\neq$ Émotion :} Dans CyberAggAdo, l'agressivité n'est pas toujours émotionnelle. Le modèle projette souvent l'agressivité vers \textit{Colère} ou \textit{Autre}.
        \item \textbf{Le Dégoût est invisible :} Les probabilités plafonnent très bas. Le dégoût est exprimé par le rejet social ou la dévalorisation corporelle, des formes absentes de TextToKids.
    \end{itemize}
\end{frame}

\begin{frame}{Pourquoi le modèle se trompe ? (2/2)}
    \begin{itemize}
        \item \textbf{Les modes implicites s'effondrent :} \textit{Suggérée} et \textit{Comportementale} exigent de comprendre une situation sociale (menace, humiliation, exclusion).
        \item \textbf{Bruit du contexte BCA :} Conçu pour textes continus. Dans CyberAggAdo, le contexte multi-locuteurs contamine la phrase cible.
        \item \textbf{Bruit lexical et Argot :} Les abréviations (\textit{tg, frr, ptn}) portent un sens pragmatique fort mais sont mal gérées par le tokenizer de CamemBERT.
    \end{itemize}
\end{frame}

\begin{frame}{Synthèse : Les causes principales}
    \begin{enumerate}
        \item \textbf{Domain shift textuel \& pragmatique} (actes sociaux de harcèlement).
        \item \textbf{Labels rares ignorés} (\textit{Dégoût, Culpabilité}).
        \item \textbf{Confusion agression / émotion}.
        \item \textbf{Modes d'expression non transférables}.
        \item \textbf{Bruit lexical} mal représenté par le tokenizer.
        \item \textbf{Contexte inadapté} aux chats multi-locuteurs.
        \item \textbf{Indépendance des têtes} générant des incohérences.
    \end{enumerate}
\end{frame}

\begin{frame}{Recommandations}
    \begin{itemize}
        \item \textbf{Calibration des seuils} : Fixer des seuils spécifiques (ex: abaisser pour le Dégoût), bien que cela soit insuffisant seul.
        \item \textbf{Post-traitement logique} : Forcer \textit{Emo=1} si une catégorie est prédite, attribuer un mode si manquant, etc.
        \item \textbf{Évaluer avec et sans contexte} : Éviter le contexte sur des échantillons non contigus.
        \item \textbf{Normalisation lexicale} : Dictionnaire d'abréviations et contrôle des élongations.
        \item \textbf{Adaptation de domaine / Fine-tuning} sur CyberAggAdo essentiel pour les catégories complexes.
    \end{itemize}
\end{frame}

\begin{frame}{Enrichissement : Analyse d'erreurs (SHAP \& FP-Growth)}
    \begin{itemize}
        \item L'erreur augmente significativement avec la \textbf{densité émotionnelle}, la longueur du texte, le \textbf{mépris/haine}, et l'agression ouverte (\texttt{HATE=OAG}).
        \item \textbf{SHAP} : Confirme l'importance de \texttt{text\_length}, \texttt{mépris/haine}, et de la cible interactionnelle (\texttt{TARGET}).
        \item \textbf{MLxtend} : Les règles d'association montrent que les formes non standard (abréviations, élongations) associées à l'agressivité concentrent les échecs.
        \item L'échec d'EMOTYC est d'abord dû à l'inadéquation pragmatique du domaine.
    \end{itemize}
\end{frame}

\begin{frame}{Conclusion}
    \begin{itemize}
        \item EMOTYC détecte la présence globale d'émotion, mais échoue sur les spécificités : catégories fines, modes, distinction agression/émotion, signaux pragmatiques.
        \item \textbf{Différence clé} : TextToKids (structures narratives/éditées) vs. CyberAggAdo (actes de parole, rôles sociaux, implicite).
        \item \textbf{Perspectives} : Calibrer les seuils, appliquer des contraintes logiques, et prévoir une adaptation spécialisée au domaine du cyberharcèlement adolescent.
    \end{itemize}
\end{frame}

\end{document}
